\subsection{凸性定理}

定理 1. --- 设 X 是一个局部紧空间，\(\mu\) 是 X 上的一个正测度，F 是一个实 Banach 空间，D 是 F 中的一个闭凸集，f 是 X 上的一个函数，使得 \(\mathbf{f}(\mathbf{X}) \subset \mathbf{D}\) 。对每个非可忽略的可积数值函数 \(g \geqslant 0\)（使 \(\mathbf{fg}\) 可积者），点 \(\frac{\int \mathbf{fg} d\mu}{\int g d\mu}\) 属于 D。

因为，设 \(\mathbf{F}'\) 是 \(\mathbf{F}\) 的对偶，而 \(\langle \mathbf{z},\mathbf{a}'\rangle \leqslant \alpha (\mathbf{a}'\in \mathbf{F}',\alpha \in \mathbf{R})\) 是定义一个包含 D 的闭半空间的关系式。由于 \(\mathbf{fg}\) 可积，数值函数 \(\langle \mathbf{fg},\mathbf{a}'\rangle = \langle \mathbf{f},\mathbf{a}'\rangle g\) 也可积，且

\[
\int \left\langle \mathbf {f} g, \mathbf {a} ^ {\prime} \right\rangle d \mu = \left\langle \int \mathbf {f} g d \mu , \mathbf {a} ^ {\prime} \right\rangle
\]

（§4, No.~2, 定理 1 的推论 1）；但由假设，\(\langle \mathbf{f}(x), \mathbf{a}' \rangle \leqslant \alpha\) 对所有 \(x \in X\) 成立，因此 \(\langle \mathbf{f}(x)g(x), \mathbf{a}' \rangle \leqslant \alpha g(x)\)；积分之，得

\[
\left\langle \int \mathbf {f} g d \mu , \mathbf {a} ^ {\prime} \right\rangle \leqslant \alpha \int g d \mu .
\]

这就证明了点 \(\frac{\int\mathbf{f}g d\mu}{\int g d\mu}\) 属于包含 D 的每个闭半空间；而由 Hahn--Banach 定理，D 是包含它的诸闭半空间之交（TVS, II, §5, No.~3, 命题 4 的推论 1），定理得证。

推论. --- 若正测度 \(\mu\) 的总质量等于 1，且 \(\mathbf{f}\) 可积，则 \(\int \mathbf{f} d\mu\) 属于 \(\mathbf{f}(\mathbf{X})\) 在 \(\mathbf{F}\) 中的闭凸包。

只需对 g 取常值函数 1。

